% BEGIN REV09_MULTIPLICIDAD_MODAL
\subsubsection{Multiplicidad y recuperación de itinerarios modales}
\label{corpus9:modal:multiplicidad}

La incidencia \(F_{\mathrm{av}}\), construida por el calendario trítico
en el desarrollo precedente, actúa aquí sobre su envolvente de paseos
modales. Escribimos
\[
 F=F_{\mathrm{av}}=\begin{pmatrix}0&1\\1&1\end{pmatrix},
 \qquad 0=\Sigma,\quad 1=\Pi.
\]
La coordenada de autoescala \(\varphi\) y el rayo positivo
\(r=(1,\varphi)\) son las salidas ya obtenidas de esa incidencia.
Sobre ellas se determina la multiplicidad de la lectura escalar y
se construye una representación reversible de cada itinerario.

Se considera el conjunto de paseos dirigidos finitos
\(\gamma=(x_0,\ldots,x_n)\) de \(F\). Se distinguen sus extremos
y longitudes, se admiten repeticiones y se incluyen los paseos triviales.
Este dominio es el envolvente modal; el estado completo del TPK
conserva, además, los registros de su historia enriquecida.
El carácter multiplicativo adopta la forma
\begin{equation}
 \chi(\gamma)=\varphi^n\frac{r_{x_0}}{r_{x_n}}
             =\varphi^{k(\gamma)},\qquad
 r=(1,\varphi),\quad k(\gamma)=n+x_0-x_n.
 \label{corpus9:modal:caracter}
\end{equation}
Las transiciones \(\Sigma\to\Pi\), \(\Pi\to\Pi\) y
\(\Pi\to\Sigma\) incrementan \(k\) en \(0\), \(1\) y \(2\),
respectivamente. El primer paso modifica la ruta manteniendo su
carácter. En particular, \((1,1)\) y \((0,1,1)\) producen ambos
\(\varphi\) bajo esta lectura.

\begin{proposition}[Multiplicidad de la lectura de autoescala]
\label{corpus9:modal:conteo}
Sea \(\mathcal F_k=\{\gamma:k(\gamma)=k\}\) y
\(N_k=|\mathcal F_k|\). Entonces
\begin{equation}
 N_0=3,\qquad
 N_k=2\operatorname{tr}(F^k)
     =2\bigl(\varphi^k+(-\varphi^{-1})^k\bigr)
     \quad(k\geq1),
 \label{corpus9:modal:multiplicidad-exacta}
\end{equation}
y
\[
 \lim_{k\to\infty}N_k^{1/k}=\varphi.
\]
\end{proposition}

\begin{proof}
Dados los extremos \(i,j\in\{0,1\}\), la longitud está determinada
por \(n=k-i+j\). El número de paseos con esos extremos es
\((F^n)_{ij}\). Para \(k\geq1\), la suma de los cuatro sectores da
\[
 N_k=\operatorname{tr}(F^k)
      +(F^{k+1})_{01}+(F^{k-1})_{10}.
\]
La identidad \(F^2=F+I\) implica, por inducción, que para \(n\geq1\)
\[
 F^n=\begin{pmatrix}f_{n-1}&f_n\\f_n&f_{n+1}\end{pmatrix},
 \qquad f_0=0,\quad f_1=1,\quad f_{n+1}=f_n+f_{n-1}.
\]
Para \(k\geq2\), los dos términos no diagonales son
\(f_{k+1}+f_{k-1}=\operatorname{tr}(F^k)\); para \(k=1\),
su suma es \(1+0=1=\operatorname{tr}(F)\), usando \(F^0=I\).
Para \(k=0\), los únicos paseos son \((0)\), \((1)\) y \((0,1)\).
Los autovalores de \(F\) son \(\varphi\) y \(-\varphi^{-1}\),
por lo que su traza proporciona la fórmula cerrada.
Finalmente,
\[
 N_k=2\varphi^k\bigl(1+(-\varphi^{-2})^k\bigr),
\]
y el factor entre paréntesis tiende a uno; tomando raíces
\(k\)-ésimas se obtiene el límite.
\end{proof}

La autoescala coincide así con la tasa exponencial de crecimiento
de las genealogías modales finitas que el carácter escalar
hace indistinguibles.

\paragraph{Refinamiento por composición de los retornos.}
Sean \(b(\gamma)\) el número de aristas \(\Pi\to\Sigma\) y
\(c(\gamma)\) el de aristas \(\Pi\to\Pi\). Los incrementos de
\eqref{corpus9:modal:caracter} dan \(k=2b+c\).
Con \(z,t\) indeterminadas formales, la incidencia graduada
\[
 M(z,t)=\begin{pmatrix}0&1\\tz^2&z\end{pmatrix}
\]
registra simultáneamente exponente y retornos. La única arista
de grado cero es \(\Sigma\to\Pi\), y carece de ciclo de grado cero.
Por ello cada coeficiente cuenta un conjunto finito de paseos.
Si \(N_{k,b}\) denota su multiplicidad con valores \(k,b\), la suma
de todas las longitudes satisface
\begin{equation}
 \sum_{k,b\geq0}N_{k,b}z^kt^b
 =\mathbf1^{\mathsf T}(I-M(z,t))^{-1}\mathbf1
 =\frac{3-z+tz^2}{1-z-tz^2},
 \qquad \mathbf1=(1,1)^{\mathsf T}.
 \label{corpus9:modal:serie-graduada}
\end{equation}
En efecto, la entrada \(ij\) de \(M^n\) cuenta los paseos de
longitud \(n\) con esos extremos y sus pesos. La suma formal
coeficiente a coeficiente proporciona \((I-M)^{-1}\).
El determinante de \(I-M\) es \(1-z-tz^2\), y la suma de las
cuatro entradas de su adjunta es \(3-z+tz^2\).

Para \(k\geq1\) y \(0\leq b\leq\lfloor k/2\rfloor\), resulta
\begin{equation}
 N_{k,b}=\frac{2k}{k-b}\binom{k-b}{b}.
 \label{corpus9:modal:conteo-retornos}
\end{equation}
Para demostrarlo, la expansión
\[
 (1-z-tz^2)^{-1}=\sum_{m\geq0}(z+tz^2)^m
\]
da el coeficiente \(\binom{k-b}{b}\) en \(z^kt^b\).
La multiplicación por el numerador de
\eqref{corpus9:modal:serie-graduada} produce
\[
 3\binom{k-b}{b}-\binom{k-b-1}{b}
   +\binom{k-b-1}{b-1}
 =2\binom{k-b}{b}+2\binom{k-b-1}{b-1}.
\]
Se toman como cero los coeficientes fuera de rango.
Para \(b\geq1\), la igualdad
\(\binom{k-b-1}{b-1}=\frac{b}{k-b}\binom{k-b}{b}\)
da \eqref{corpus9:modal:conteo-retornos}; para \(b=0\), ambos
miembros valen dos. El término constante se separa:
\(N_{0,0}=3\). Para \(k=4\) hay catorce paseos, distribuidos en
\(2\), \(8\) y \(4\) según contengan cero, uno o dos retornos.
Por tanto, el mismo carácter admite composiciones internas distintas.

\begin{proposition}[Recuperación reversible y transporte de códigos]
\label{corpus9:modal:recuperador}
Cada paseo del envolvente modal admite un código
\[
 \mathcal C(\gamma)=(k,a),\qquad 0\leq a<N_k,
\]
con decodificación \(\mathcal D\) tal que
\[
 \mathcal D\mathcal C=\mathrm{id},\qquad
 \mathcal C\mathcal D=\mathrm{id}.
\]
Si \(A_v(\gamma)=\gamma v\) prolonga por una arista admisible y
\(T\) elimina el último estado de un paseo de longitud positiva,
sus representaciones
\[
 \mathcal E_v=\mathcal C A_v\mathcal D,\qquad
 \widehat T=\mathcal C T\mathcal D
\]
satisfacen, en sus dominios respectivos,
\begin{equation}
 \widehat T\mathcal E_v=\mathrm{id},\qquad
 k(A_v\gamma)-k(\gamma)=1+x_n-v.
 \label{corpus9:modal:transporte-codigos}
\end{equation}
\end{proposition}

\begin{proof}
Para cada \(k\), se ordenan los sectores no vacíos por \((n,i,j)\)
y, dentro de cada sector, los paseos lexicográficamente por sus
estados. El número de continuaciones de longitud \(m\) desde
\(v\) hasta \(j\) es \((F^m)_{vj}\). Al recorrer un paseo se suman
los tamaños de los sectores precedentes y, en cada posición, los
de las continuaciones lexicográficamente anteriores a la elegida.
El resultado es su índice \(a\), contado desde cero.

Para invertir el proceso, se resta del índice el tamaño de cada
sector precedente hasta localizar el único sector que lo contiene.
En cada posición se procede del mismo modo con los tamaños de
las continuaciones. Estos conjuntos forman una partición disjunta
y exhaustiva del sector de partida. Para longitud cero, el único
paseo de un sector no vacío es su vértice. Suponiendo recuperadas
las continuaciones de longitud \(m-1\), los intervalos de índices
de cada primer sucesor identifican unívocamente ese sucesor y el
índice restante de su continuación. La inducción en \(m\) prueba
que ambos procedimientos son inversos, tanto sobre paseos como
sobre los índices \(0,\ldots,N_k-1\).

El orden lexicográfico serializa la ruta existente; las operaciones
que generan la ruta preceden a esa serialización.
Como \(T A_v=\mathrm{id}\), las dos identidades inversas dan
\[
 \widehat T\mathcal E_v
 =\mathcal C T\mathcal D\mathcal C A_v\mathcal D
 =\mathcal C T A_v\mathcal D
 =\mathrm{id}.
\]
Finalmente, la definición del exponente da
\[
 k(A_v\gamma)=(n+1)+x_0-v,\qquad
 k(\gamma)=n+x_0-x_n,
\]
cuya diferencia es la segunda identidad.
\end{proof}

\paragraph{Capacidad de distinción con exponente conocido.}
Conocido \(k\), una representación binaria inyectiva de longitud
fija \(\ell\) requiere \(2^\ell\geq N_k\), por el principio de
las casillas. En consecuencia,
\(\ell\geq\lceil\log_2N_k\rceil\).
El índice \(a\) anterior alcanza esa longitud mediante su
representación binaria, completada con ceros iniciales cuando sea
necesario. Este mínimo corresponde a la transmisión del índice
con \(k\) conocido; la transmisión de \(k\) y la delimitación de
códigos de longitud variable tienen costes propios. Además,
\[
 \lim_{k\to\infty}\frac{\log_2N_k}{k}=\log_2\varphi.
\]
La igualdad mide capacidad combinatoria de distinción de rutas
en el envolvente modal.

\paragraph{Información de fibra bajo una distribución explícita.}
Fijado \(k\), sea \(R_k\) un paseo elegido uniformemente en
\(\mathcal F_k\). Usamos logaritmos naturales y
\(H(R_k)=-\sum_{\gamma\in\mathcal F_k}p_\gamma\log p_\gamma\),
con \(p_\gamma=1/N_k\). Todos estos paseos tienen carácter
\(\varphi^k\), de modo que
\[
 H(R_k\mid\chi(R_k))=\log N_k,\qquad
 H(R_k\mid\mathcal C(R_k))=0.
\]
La primera igualdad procede de que la lectura es constante en la
fibra; la segunda, de la decodificación única ya demostrada.
En particular, \(\log N_k/k\to\log\varphi\).

Para \(k=4\), el carácter deja una incertidumbre de \(\log14\)
nats. Las tres clases \(B=b(R_4)\) tienen probabilidades
\(2/14\), \(8/14\) y \(4/14\). Condicionada a cada clase, la
distribución permanece uniforme en sus \(2\), \(8\) o \(4\)
elementos, respectivamente. Por definición de entropía
condicional,
\[
 H(R_4\mid B)
 =\frac2{14}\log2+\frac8{14}\log8+\frac4{14}\log4.
\]
El índice dentro de la clase, junto con \(B\), recupera unívocamente
el paseo. Estos valores pertenecen a la distribución uniforme
de la fibra modal declarada. La medida del espacio completo de
historias, la realización temporal de \(k\) y cualquier lectura
térmica conservan sus dominios y transductores respectivos.
% END REV09_MULTIPLICIDAD_MODAL
